\documentclass[11pt]{article}
\usepackage{fontspec}
\setmainfont{Times New Roman}
\setsansfont{Arial}
\setmonofont{Consolas}
\usepackage{amsmath,amssymb,mathtools}
\usepackage{geometry}
\usepackage{microtype}
\newcommand{\mS}{\mathfrak{S}}
\newcommand{\mA}{\mathfrak{A}}
\newcommand{\mat}[1]{\begin{pmatrix}#1\end{pmatrix}}
\geometry{a4paper,margin=2.35cm}
\setlength{\parindent}{1.25em}
\setlength{\parskip}{0pt}
\makeatletter
\renewcommand\footnotesize{\@setfontsize\footnotesize{7.6}{8.6}\selectfont}
\makeatother
\emergencystretch=100em
\allowdisplaybreaks
\sloppy
\hbadness=10000
\vbadness=10000
\hfuzz=2pt
\vfuzz=10pt
\raggedbottom

\begin{document}

% Unit: Paper32 section 3 cyclic quaternion fields v001. Direct Russian translation from the R122 German Paper32/Paper33 source slice, lines 482--516 / segments P32-S0023--P32-S0029.
\noindent Далее следует показать, что для бесконечно многих $n$ существуют циклические поля степени $2^n$, являющиеся минимальными полями расщепления группы кватернионов.

\medskip

\noindent Пусть $p$ -- рациональное простое число, для которого при некотором $r>0$
\[
        (*)\qquad 2^r+1\equiv0\pmod p
\]
выполняется; пусть $2^n$ есть наивысшая степень числа $2$, делящая $p-1$.

Сначала покажем, что для бесконечно многих значений $n$ существуют простые числа $p$. Для этого пусть $k$ -- произвольное положительное целое рациональное число, а $p$ -- простой делитель числа $2^{2^k}+1$. Тогда требуемая для $p$ конгруэнция выполнена при $r=2^k$. Кроме того, $2\pmod p$ имеет порядок $2^{k+1}$, и поэтому для наивысшей степени $2^n$ числа $2$, делящей $p-1$, число $n$ больше $k$. Во всяком случае, значит, существуют сколь угодно большие $n$, которым соответствуют простые числа $p$.

Пусть $\varepsilon$ -- первообразный $p$-й корень из единицы. Теперь покажем, что в $P(\varepsilon)$ число $-1$ представляется как сумма двух квадратов
\[
        -1=a^2+b^2=(a+ib)(a-ib)\qquad (a,b\text{ числа из }P(\varepsilon)).
\]
Это равносильно тому, что в поле $(4p)$-х корней из единицы $P(\varepsilon,i)=P(\varepsilon\cdot i)$ существуют числа, относительная норма которых относительно $P(\varepsilon)$ равна как раз $-1$.

Число $1+i\varepsilon^\nu$ $(\nu=1,2,\ldots,p-1)$ имеет относительную норму $(1+i\varepsilon^\nu)(1-i\varepsilon^\nu)=1+\varepsilon^{2\nu}$; следовательно, все числа $1+\varepsilon^\nu$ $(\nu=1,2,\ldots,p-1)$ являются относительными нормами чисел из $P(\varepsilon i)$ относительно $P(\varepsilon)$ как основного поля. Поэтому то же самое верно и для
\[
\Pi=\prod_{\nu=0}^{r-1}(1+\varepsilon^{2^\nu})
=(1+\varepsilon)(1+\varepsilon^2)(1+\varepsilon^4)\cdots(1+\varepsilon^{2^{r-1}})
=\sum_{\lambda=0}^{2^r-1}\varepsilon^\lambda.
\]
Если к $\Pi$ добавить еще $\varepsilon^{2^r}=\varepsilon^{-1}=\varepsilon^{p-1}$ (в силу $(*)$), то $\Pi+\varepsilon^{p-1}$ становится суммой целых частей выражения $1+\varepsilon+\varepsilon^2+\cdots+\varepsilon^{p-1}=0$, и потому
\[
        \Pi=-\varepsilon^{-1}.
\]
Теперь и $\varepsilon$ является относительной нормой числа из $P(\varepsilon i)$, а именно числа $\varepsilon^{(p+1)/2}$. Следовательно, аналогичное верно для
\[
        \Pi\varepsilon=-1.
\]
Поэтому $-1$ в $P(\varepsilon)$ представим как сумма двух квадратов;\footnote{Для простых чисел вида $8n\pm3$ это рассуждение уже имеется у E. Landau, Über die Darstellung definiter Funktionen durch Quadrate, Math. Ann. Bd. 62, с. 272--285.} значит, $P(\varepsilon)$ является полем расщепления для группы кватернионов.

Пусть $K$ -- циклическое поле степени $2^n$, содержащееся в $P(\varepsilon)$; мы утверждаем, что и $K$ является полем расщепления.\footnote{Выбор циклического поля расщепления степени $2^n$ как подполя поля деления круга повторяет пример Хассе.} Пусть $m$ -- индекс тела кватернионов относительно $K$ как основного поля; тогда $m=1$ или $m=2$. Но поскольку существует поле расщепления $P(\varepsilon)$ нечетной относительной степени $(p-1)/2^n$ относительно $K$, второй случай невозможен; значит, $m=1$. Это и означает, что $K$ является полем расщепления.

Итак, для бесконечно многих $n$ существуют циклические поля степени $2^n$, являющиеся минимальными полями расщепления.

\end{document}
